
\documentclass{article}
\usepackage{colortbl}
\usepackage{makecell}
\usepackage{multirow}
\usepackage{supertabular}

\begin{document}

\newcounter{utterance}

\centering \large Interaction Transcript for game `cladder', experiment `full\_v1.5\_default', episode 8207 with Bentel rockers\_\_qwen3.5{-}4b\_\_Bentel\_iter\_3.
\vspace{24pt}

{ \footnotesize  \setcounter{utterance}{1}
\setlength{\tabcolsep}{0pt}
\begin{supertabular}{c@{$\;$}|p{.15\linewidth}@{}p{.15\linewidth}p{.15\linewidth}p{.15\linewidth}p{.15\linewidth}p{.15\linewidth}}
    \# & \multicolumn{2}{c}{Player} && \multicolumn{2}{c}{Game Master} \\
    \hline

    \theutterance \stepcounter{utterance}  
    & & & \multicolumn{4}{p{0.6\linewidth}}{
        \cellcolor[rgb]{0.9,0.9,0.9}{
            \makecell[{{p{\linewidth}}}]{
                \texttt{\tiny{[P1$\langle$GM]}}
                \texttt{You are an expert in causal inference. The following question is not a typical commonsense query, but rather a meticulously designed question created by a professor specializing in causal inference, intended to assess the students' mastery of the course content.} \\
\texttt{Imagine a self{-}contained, hypothetical world with only the following conditions, and without any unmentioned factors or causal relationships: Tijv has a direct effect on tijw and xevo. Xevo has a direct effect on gyzp. Tijw has a direct effect on gyzp. For those who are not tijv and are not xevo, the probability of gyzp is 29\%. For those who are not tijv and are xevo, the probability of gyzp is 14\%. For those who are tijv and are not xevo, the probability of gyzp is 11\%. For those who are tijv and are xevo, the probability of gyzp is 16\%. The overall probability of tijv is 88\%. Will xevo increase the chance of gyzp?} \\
\texttt{Start your answer with 'Yes' or 'No', followed by additional reasoning or evidence to support your explanation.} \\
            }
        }
    }
     \\ \\

    \theutterance \stepcounter{utterance}  
    & \multicolumn{4}{p{0.6\linewidth}}{
        \cellcolor[rgb]{0.561,0.737,0.561}{
            \makecell[{{p{\linewidth}}}]{
                \texttt{\tiny{[P1$\rangle$GM]}}
                \texttt{No} \\
\\ 
\texttt{**Reasoning:**} \\
\\ 
\texttt{To determine whether **xevo** increases the chance of **gyzp**, we must compare the probability of **gyzp** conditional on **xevo** being present versus absent, while holding the status of the other causal factor, **tijv**, constant. This is necessary because **xevo** and **tijv** are both direct causes of **gyzp**, creating a confounding structure where the effect of one variable depends on the state of the other.} \\
\\ 
\texttt{We are given the following conditional probabilities for **gyzp**:} \\
\texttt{1.  \$P(\textbackslash{}text\{gyzp\} \textbackslash{}mid \textbackslash{}neg \textbackslash{}text\{tijv\}, \textbackslash{}neg \textbackslash{}text\{xevo\}) = 29\textbackslash{}\%\$} \\
\texttt{2.  \$P(\textbackslash{}text\{gyzp\} \textbackslash{}mid \textbackslash{}neg \textbackslash{}text\{tijv\}, \textbackslash{}text\{xevo\}) = 14\textbackslash{}\%\$} \\
\texttt{3.  \$P(\textbackslash{}text\{gyzp\} \textbackslash{}mid \textbackslash{}text\{tijv\}, \textbackslash{}neg \textbackslash{}text\{xevo\}) = 11\textbackslash{}\%\$} \\
\texttt{4.  \$P(\textbackslash{}text\{gyzp\} \textbackslash{}mid \textbackslash{}text\{tijv\}, \textbackslash{}text\{xevo\}) = 16\textbackslash{}\%\$} \\
\\ 
\texttt{We also know the marginal probability of **tijv** is \$88\textbackslash{}\%\$, so \$P(\textbackslash{}neg \textbackslash{}text\{tijv\}) = 12\textbackslash{}\%\$. However, the marginal probability is not required to answer the specific causal question regarding the direction of the effect of **xevo** on **gyzp** within the defined world.} \\
\\ 
\texttt{**Step 1: Analyze the effect of xevo when tijv is absent.**} \\
\texttt{*   If a subject is **not** **tijv** and **is** **xevo**: Probability = \$14\textbackslash{}\%\$} \\
\texttt{*   If a subject is **not** **tijv** and **is not** **xevo**: Probability = \$29\textbackslash{}\%\$} \\
\texttt{*   Comparison: \$14\textbackslash{}\% < 29\textbackslash{}\%\$. In this stratum, **xevo** *decreases* the probability of **gyzp**.} \\
\\ 
\texttt{**Step 2: Analyze the effect of xevo when tijv is present.**} \\
\texttt{*   If a subject is **tijv** and **is** **xevo**: Probability = \$16\textbackslash{}\%\$} \\
\texttt{*   If a subject is **tijv** and **is not** **xevo**: Probability = \$11\textbackslash{}\%\$} \\
\texttt{*   Comparison: \$16\textbackslash{}\% > 11\textbackslash{}\%\$. In this stratum, **xevo** *increases* the probability of **gyzp**.} \\
\\ 
\texttt{**Conclusion:**} \\
\texttt{The effect of **xevo** on **gyzp** is not uniform; it depends entirely on the presence or absence of **tijv**.} \\
\texttt{*   When **tijv** is present, **xevo** increases the chance of **gyzp** (by 5 percentage points).} \\
\texttt{*   When **tijv** is absent, **xevo** decreases the chance of **gyzp** (by 15 percentage points).} \\
\\ 
\texttt{In causal inference, a variable is said to "increase the chance" of an outcome only if its causal effect is positive across the relevant population or under specific counterfactual conditions. Since **xevo** has a negative causal effect in the majority of the population (where **tijv** is absent, which constitutes 12\% of the population, but more importantly, represents a distinct causal pathway where the sign of the effect flips), we cannot universally state that **xevo** increases the chance of **gyzp**. In fact, for the subset of individuals without **tijv**, it actively reduces the probability. Therefore, the statement "xevo will increase the chance of gyzp" is false as a general rule in this hypothetical world.} \\
\\ 
\texttt{*(Note: If the question implies a marginal comparison \$P(\textbackslash{}text\{gyzp\} \textbackslash{}mid \textbackslash{}text\{xevo\})\$ vs \$P(\textbackslash{}text\{gyzp\} \textbackslash{}mid \textbackslash{}neg \textbackslash{}text\{xevo\})\$, the answer would still be No, because the negative effect in the \$\textbackslash{}neg \textbackslash{}text\{tijv\}\$ group outweighs the positive effect in the \$\textbackslash{}text\{tijv\}\$ group given the prevalence of \$\textbackslash{}neg \textbackslash{}text\{tijv\}\$, leading to a net decrease. However, the stratum{-}specific analysis above provides the definitive causal mechanism showing the reversal of effects.)*} \\
            }
        }
    }
    & & \\ \\

    \theutterance \stepcounter{utterance}  
    & & & \multicolumn{2}{p{0.3\linewidth}}{
        \cellcolor[rgb]{0.9,0.9,0.9}{
            \makecell[{{p{\linewidth}}}]{
                \texttt{\tiny{[GM$|$GM]}}
                \texttt{No} \\
            }
        }
    }
    & & \\ \\

    \theutterance \stepcounter{utterance}  
    & & & \multicolumn{2}{p{0.3\linewidth}}{
        \cellcolor[rgb]{0.9,0.9,0.9}{
            \makecell[{{p{\linewidth}}}]{
                \texttt{\tiny{[GM$|$GM]}}
                \texttt{yes} \\
            }
        }
    }
    & & \\ \\

    \theutterance \stepcounter{utterance}  
    & & & \multicolumn{2}{p{0.3\linewidth}}{
        \cellcolor[rgb]{0.9,0.9,0.9}{
            \makecell[{{p{\linewidth}}}]{
                \texttt{\tiny{[GM$|$GM]}}
                \texttt{game\_result = LOSE} \\
            }
        }
    }
    & & \\ \\

\end{supertabular}
}

\end{document}
